\section{Background: optimization in energy system modeling}
\label{sec:background}

Surrogate models are useful only insofar as they sit inside an
optimization problem; the choice of surrogate family is therefore not
independent of the optimization formulation it has to live with. This
section sets the optimization-side vocabulary that the rest of the
review uses and identifies, for each problem class, where the
computational pressure that motivates surrogate use actually comes
from.

% =====================================================================
% Subsection 2.1 - Problem classes
% =====================================================================

\subsection{Problem classes}
\label{sec:background:problem-classes}

\paragraph{Linear and mixed-integer dispatch.}
The largest single class of energy-system optimization problems is
short-term operational dispatch in linear (LP) or mixed-integer linear
(MILP) form. Economic dispatch, security-constrained economic dispatch
and unit commitment are formulated as MILPs whenever start-up costs,
minimum-up/down times or block start-ups are modelled
\cite{Wu2022231,Chen20244723,Chen2022,Lu20232989,Safta20172535,Pang2023,Wang20211767,Lin2023,Dullinger20181646,Nikolaidis2021}.
The same MILP backbone underpins the operational layer of capacity
expansion and multi-energy planning models
\cite{Rodgers2018951,Jahangiri20244849,Jahangiri2023,Han2018622,Fu20202904,Fu2022,Cao2023,Yin2024,Wang2020202933,Ren2024}.
On the network side, optimal power flow (OPF) in either DC or AC
form, and security- and chance-constrained variants thereof, define
the rest of the dispatch landscape
\cite{Muhlpfordt20194806,Su2022315,Deng2020831,Liu20223726,Chen2023657,Koirala20242284,Xu20212696,Srithapon202134395,Urgun20202791,Khaloie2025}.

\paragraph{Mixed-integer non-linear programming.}
Many energy applications are inherently non-linear: AC power flow
with quadratic losses, district-heating dispatch with mixed-integer
storage and flow direction, integrated electricity-heat systems with
variable mass flow, hydrogen and synthesis chains, and integrated
multi-energy hubs that include conversion characteristics
\cite{Zheng20231907,Wu2021,Cao2023,Yin2024,Ren2024,Jalving2023,Dong2023,Wang2020202933}.
Solving such MINLPs to global optimality is rarely possible at
realistic system size, which is why surrogates that linearise or
convexify the non-linear part play an outsize role in this class
\cite{Liu20223726,Chen2023657,Su2022315,Deng2020831,Park2019}.

\paragraph{Bilevel and equilibrium problems.}
Strategic-bidding, market-design and regulator-versus-operator
studies are typically posed as bilevel optimization problems with an
outer decision (regulation, investment, bidding) and an inner
clearing problem
\cite{Han2018622,Wu2022231,Wang20211767,Cao2023}.
The inner problem here is the natural target for a surrogate: it is
re-solved many times during the outer search.

\paragraph{Stochastic and robust optimization.}
Where uncertainty cannot be hedged by reasonable conservatism in a
deterministic model, two formulations dominate. Two-stage and
multi-stage stochastic programming explicitly enumerates scenarios
and computes recourse decisions
\cite{Wang20211767,Safta20172535,Hu2021671,Hu2020,Volodina2022,Lawson201647,Lin2023,Pang2023,Zhou20181}.
Robust and distributionally-robust optimization replaces the
expectation by a worst-case bound over an uncertainty set
\cite{Zhou20181,Yang2024,Dong2022,Dong2023,Liang20246508,Giorgetti2020,Fu20202904,Fu2022}.
Chance-constrained formulations sit between the two and are often
the entry point for polynomial-chaos and Gaussian-process surrogates
\cite{Muhlpfordt20194806,Koirala20242284,Xu20212696,Chen2023657,Liang20246508}.

\paragraph{Multi-objective optimization.}
Energy system design problems almost always trade off cost,
emissions, reliability, comfort and resource use, and the resulting
multi-objective formulations have driven a large share of the
surrogate-assisted-evolutionary-optimization literature
\cite{Thrampoulidis2021,Hussien20211883,Wang2017,Schulte2020,Han2018622,Roth2019214,Giorgetti2020,Hirvonen2017323,Sun2020}.
Because each evaluation along a Pareto front is itself an
optimization run, the surrogate role here is acceleration of the
outer search.

\paragraph{Model predictive and rolling-horizon control.}
A fifth class is online: model predictive control (MPC) and
rolling-horizon dispatch repeatedly resolve a short-horizon problem
under updated forecasts, which means total compute is dominated by
the number of solves, not by their individual size
\cite{Dullinger20181646,Lu20232989,Faraji2020157284,Reynolds2017704,Wang20192635,Yang2024,Dong2022}.
Surrogates here are typically used to predict the recurrent solution
or to replace expensive physics that lives inside the controller.

% =====================================================================
% Subsection 2.2 - Cost drivers
% =====================================================================

\subsection{Where computational load actually comes from}
\label{sec:background:cost-drivers}

It is useful to be explicit about which problem features make energy
system optimization expensive, because the surrogate that helps for
one driver does not necessarily help for another.

\paragraph{Integer decisions.}
Start-up/shut-down decisions, block-start configurations and discrete
capacity sizes turn polynomial LP solves into exponential-worst-case
MILPs and dominate solution time once the time horizon and the unit
count grow
\cite{Wu2022231,Chen20244723,Chen2022,Pang2023,Lin2023,Lu20232989,Wang20211767}.

\paragraph{Time resolution and horizon.}
Sub-hourly resolution is often necessary to capture renewable and
flexibility dynamics; combined with multi-year horizons in
expansion planning, this drives the temporal index size to numbers
where direct solves are no longer viable
\cite{Han2018622,Fu20202904,Fu2022,Rodgers2018951,Jahangiri20244849,Jahangiri2023,Hu2021671,Volodina2022,Schulte2016104}.

\paragraph{Network detail.}
AC network constraints (voltages, reactive power, line losses) are
non-linear, and security-constrained variants multiply the
constraint set by the contingency list
\cite{Muhlpfordt20194806,Su2022315,Deng2020831,Liu20223726,Urgun20202791,Chen2023657,Koirala20242284,Xu20212696,Khaloie2025}.

\paragraph{Multi-energy coupling.}
Coupling electricity to heat, gas, hydrogen, transport and industrial
processes through conversion devices and shared resources increases
both the number of variables and the number of non-linear
relationships, and creates the integrated electricity-heat and
multi-energy planning problems for which surrogates are now common
\cite{Zhou20181,Wang2021,Zheng20231907,Ren2024,Jalving2023,Dong2023,Yin2024,Cao2023,Wang2020202933,Roth2019214}.

\paragraph{Scenario count.}
Stochastic and robust formulations inflate the model with copies of
the operational subproblem, one per scenario, and the marginal cost
of an additional scenario is high
\cite{Wang20211767,Safta20172535,Hu2021671,Hu2020,Volodina2022,Lawson201647,Liang20246508,Lin2023}.

\paragraph{Outer-loop calls.}
Multi-objective evolutionary search, Bayesian optimization, design
exploration and bilevel formulations turn a single solve into an
outer loop with hundreds to thousands of solves
\cite{Thrampoulidis2021,Hussien20211883,Wang2017,Schulte2020,Han2018622,Roth2019214,Giorgetti2020,Hirvonen2017323,Pang2023,Xiao2018}.

% =====================================================================
% Subsection 2.3 - Decomposition + surrogate hooks
% =====================================================================

\subsection{Decomposition and what surrogates can replace}
\label{sec:background:decomposition}

For each cost driver above, the optimization community has developed
decomposition strategies that surrogates can plug into
\cite{Ruan2021221,Khaloie2025,Lim2025,Wu2022231,Chen2022,Chen20244723,Liu20223726}.
Benders and L-shaped decomposition split the problem into a master
and a recourse problem; the recourse function is a natural target for
a Gaussian-process or neural surrogate that is queried millions of
times during the outer iterations
\cite{Hu2021671,Hu2020,Volodina2022,Lawson201647,Safta20172535,Liang20246508}.
Lagrangean and column-generation methods iterate between a pricing
subproblem and a master problem, and surrogate pricing or surrogate
column quality predictors have been used in recent
learning-to-optimize work
\cite{Chen20244723,Chen2022,Wu2022231,Guo20222057,Xu20223066,Li202359995}.
Scenario reduction and time-period aggregation reduce the size of
stochastic and capacity-expansion problems, and surrogates can guide
both the choice of representative scenarios and the assessment of
their loss-of-fidelity penalty
\cite{Hu2020,Volodina2022,Lin2023,Pang2023,Yang2024,Dong2023,Liang20246508,Fu20202904,Fu2022}.
At the inner level, individual subproblems -- AC power flow, building
or thermal-storage dynamics, an electrolyser conversion curve -- can
be replaced by a learned surrogate that is much cheaper per call but
still sufficient for the outer optimizer to converge
\cite{Su2022315,Deng2020831,Liu20223726,Chen2023657,Koirala20242284,Xu20212696,Park2019,Thrampoulidis2021,Han2018622,Wang2021,Roth2019214,Hirvonen2017323,Reynolds2017704}.

The recurring lesson from these case studies is that surrogates
rarely replace the entire optimizer. They replace expensive
\emph{components} -- recourse functions, AC subproblems, building
detail, scenario response surfaces -- and the value they deliver
depends on whether those components were actually the binding
computational constraint of the model
\cite{Ruan2021221,Khaloie2025,Lim2025,Mylonopoulos202332697,Starke2025214}.
Section~\ref{sec:integration} formalises this observation as the
five-pattern integration framework that organises the rest of the
review.
